\documentclass[11pt]{article}
\usepackage[margin=1in]{geometry}
\usepackage{amsmath,amssymb}
\setlength{\parindent}{0pt}
\setlength{\parskip}{6pt}
\begin{document}
\section*{STAT 412 QZ08 --- Question 9}

\subsection*{Solution}

\[
H_0:\mu_1=\mu_2,\qquad H_1:\mu_1\ne\mu_2.
\]

Using Welch's test,
\[
\text{SE}=\sqrt{\frac{0.91^2}{35}+\frac{0.56^2}{39}}=0.1780,
\]
\[
t=\frac{2.32-2.62}{0.1780}=-1.68.
\]
The Welch degrees of freedom are
\[
df\approx55.32.
\]
\[
p=0.098.
\]

Since \(p>0.01\), fail to reject \(H_0\).

The 99\% confidence interval for \(\mu_1-\mu_2\) is
\[
(\bar{x}_1-\bar{x}_2)\pm t^*SE.
\]
Using \(t^*\approx2.669\) with \(df\approx55.32\),
\[
-0.30\pm 2.669(0.1780)=(-0.775,0.175).
\]
\[
-0.775<\mu_1-\mu_2<0.175.
\]
The interval contains 0, so it supports the hypothesis-test conclusion.

\subsection*{Final Answer}

\fbox{\begin{minipage}{0.94\linewidth}
\(t=-1.68\), \(p=0.098\). Fail to reject \(H_0\). CI blanks: lower \(-0.775\), upper \(0.175\). 99\% CI: \((-0.775,0.175)\).
\end{minipage}}

\end{document}
